% ============================================================
% CAPITOLUL 8 — CONCLUZII ȘI DIRECȚII VIITOARE
% ============================================================

\chapter{Concluzii și Direcții de Dezvoltare Viitoare}
\label{ch:concluzii}

\section{Sinteza lucrării}
\label{sec:sinteza_contributii}

Lucrarea a fost reorientată în jurul modelului SIDARTH, deoarece acesta descrie
mai bine situațiile în care diagnosticul, simptomatologia și severitatea clinică
influențează evoluția observată a unei epidemii. Modelul SIR rămâne prezent, dar
rolul său este secundar: el oferă un reper matematic și numeric simplu, util pentru
comparații și verificări de regresie.

Pe această bază a fost construit un sistem software în care datele, simulările și
metricile de performanță sunt tratate unitar. API-ul Go oferă autentificare,
încărcare de dataseturi și pornire de simulări; Redis decuplează cererile
utilizatorului de execuția numerică; runner-ele R, Julia și Python rezolvă modelul
cerut prin \texttt{model\_type}; interfața Vue.js permite urmărirea, vizualizarea
și compararea rulărilor. Contribuția practică a sistemului constă în această
legătură între model, execuție și observabilitatea rezultatelor.

\section{Concluzii privind modelele}
\label{sec:concluzii_modele}

SIDARTH introduce o structură mai apropiată de realitatea raportării
epidemiologice decât SIR. Separarea cazurilor nedepistate de cele diagnosticate și
a formelor ușoare de cele severe permite formularea unor întrebări care nu pot fi
abordate natural într-un model cu trei compartimente. În același timp, această
granularitate crește numărul de parametri și face calibrarea mai dificilă.

SIR rămâne util tocmai prin simplitate. El permite verificarea conservării
populației, a pozitivității soluției și a comportamentului solverelor într-un cadru
ușor de controlat. În sistemul implementat, SIR este tratat ca model secundar, nu
ca explicație completă a datelor reale.

\section{Concluzii privind implementarea}
\label{sec:concluzii_implementare}

Implementarea arată că același contract de simulare poate fi rulat în ecosisteme
diferite fără a schimba API-ul sau interfața web. R oferă o formulare compactă
prin \texttt{deSolve}; Python oferă integrare directă cu infrastructura și
bibliotecile științifice uzuale; Julia este potrivită pentru experimente numerice
mai ambițioase, mai ales atunci când costul de compilare este amortizat.

Schimbarea schemei rezultatelor din coloane fixe \texttt{s/i/r} într-un obiect
JSONB \texttt{compartments} este una dintre modificările importante ale pivotării.
Ea face sistemul mai flexibil și permite adăugarea unor modele noi fără migrații
ale tabelei de rezultate. Costul acestei flexibilități este o serializare mai
generală și o nevoie mai mare de validare la nivel de aplicație.

Pe latura operațională, sistemul a fost deployat efectiv pe un cluster Kubernetes
v1.31 găzduit pe o platformă cloud privată, cu trei noduri control plane și trei
workeri, YugabyteDB la factor de replicare~3 și Helm pentru ciclul de viață al
bazei de date. Această experiență confirmă portabilitatea arhitecturii și pune
în evidență câteva ajustări specifice mediilor distribuite, precum aplicarea
migrațiilor prin \texttt{psql} în locul \texttt{golang-migrate} sau adăugarea
extensiei \texttt{pgcrypto} la momentul instalării, datorită bazei PostgreSQL~11
a YugabyteDB. Aceste ajustări sunt locale și nu afectează nici codul aplicativ,
nici interfața REST.

\section{Limitări}
\label{sec:limitari}

Prima limitare este calibrarea parametrilor SIDARTH. Datele OWID folosite în
lucrare nu oferă direct distribuția pe compartimentele $D$, $A$, $R$ și $T$.
Inițializarea folosită de sistem este reproductibilă, dar nu înlocuiește o
estimare epidemiologică dedicată. Pentru rezultate interpretate clinic, modelul
trebuie calibrat pe date despre testare, spitalizare, terapie intensivă și
vindecări.

A doua limitare privește evaluarea empirică. Infrastructura de benchmark este
implementată, dar concluziile cantitative necesită rulări repetate pe configurații
hardware documentate. Fără aceste campanii complete, analiza performanței rămâne un
cadru de măsurare și nu o demonstrație statistică a superiorității unui limbaj.

O altă limitare este nivelul de maturitate al aplicației web. Fluxurile principale
sunt acoperite, însă lipsesc încă teste automate end-to-end, mecanisme de retry mai
sofisticate pentru job-uri, politici de rate limiting și audit al operațiilor
utilizatorilor. Aceste elemente ar fi necesare înainte de folosirea sistemului într-un
context operațional.

\section{Direcții viitoare}
\label{sec:directii_viitoare}

Direcția imediată este calibrarea modelului SIDARTH pe date mai bogate decât
seria de cazuri active. O extensie naturală este reintroducerea compartimentului
$E$ din SIDARTHE, pentru a modela explicit decesele, și adăugarea unor variațiuni
precum SEIR sau SIRS în același contract de simulare.

La nivel de sistem, următorul pas este completarea campaniilor de benchmark și
adăugarea unei suite automate de teste. Pentru backend, aceasta înseamnă teste
unitare și de integrare cu o bază de date izolată; pentru frontend, teste care
verifică fluxurile reale din browser. În zona de infrastructură, ar fi utilă
introducerea unor mecanisme explicite de retry și dead-letter queue pentru job-urile
care eșuează repetat, precum și activarea autoscalării orizontale (HPA) a
runner-elor pe metrici personalizate (de exemplu lungimea cozii Redis) și a
politicilor de rețea (NetworkPolicy) pentru izolare zero-trust intra-cluster.
Terminarea TLS automată prin \texttt{cert-manager} și un stack de
observabilitate (Prometheus + Grafana) ar completa profilul operațional al
deployment-ului Kubernetes.

Pe termen mai lung, sistemul poate fi extins cu vizualizări geografice, surse de
date actualizate periodic și module de alertare. Aceste extensii trebuie construite
pe aceeași disciplină metodologică: ipotezele modelului, calitatea datelor și
metricile software trebuie să rămână vizibile.
